\chapter{Retour nonadique, mémoire et holonomie de Cartan}
\label{chap:puente-retorno-holonomia-cartan}
\index{retour!27--54--108}\index{holonomie!TPK--Cartan}
\index{mémoire!distinction par rapport à la torsion}

La réalisation gravitationnelle ne commence pas par identifier une variable discrète
à une grandeur géométrique. Elle commence par distinguer trois niveaux : le retour
du transport observable, la mémoire qui différencie des histoires ayant la même
projection et l'holonomie d'une connexion. Le pont formulé ici
reçoit les retours et les cocycles déjà démontrés dans HMT ; il ne les reconstruit pas
à partir de l'action de Holst, d'une masse ou d'une constante métrologique.

\section{Hiérarchie exacte des retours}
\label{sec:jerarquia-retornos-27-54-108}

Soit \(F_{\rm obs}\) la chronologie observable définie dans le noyau HMT et soit
\[
 \mathcal K_{27}:=F_{\rm obs}^{27}.
\]
Le symbole \(\mathcal K_{27}\) désigne exclusivement la composition de
vingt-sept mises à jour de la boucle positionnelle et orientée ; il ne désigne pas une
fenêtre générique, une trace de \(108\) pas ni l'état TPK enrichi
complet. Notons \(p\) la projection sur les positions des deux
curseurs et \(\mathsf J_{\rm or}\) l'inversion simultanée de leurs
orientations. La démonstration primaire se trouve dans
\cref{prop:retorno-fobs-identidad} ; ses conséquences typées sont
\[
 p(\mathcal K_{27}x)=p(x),\qquad
 o(\mathcal K_{27}x)=\mathsf J_{\rm or}\,o(x),
 \qquad \mathsf J_{\rm or}^{\,2}=I,
\]
\[
 F_{\rm obs}^{54}:
 \text{rétablit les positions et les orientations},
 \qquad
 F_{\rm obs}^{108}=I:
 \text{rétablit aussi la phase observable}.
\]
Par conséquent, retour de position, retour d'orientation et retour de phase sont
des affirmations différentes. Aucune d'elles n'implique à elle seule le retour de
tous les champs d'une extension enrichie.

Le quotient de modes et les cocycles d'orientation sont établis dans
\cref{sec:cociclos-modo-orientacion,thm:descenso-necesario-fav}. Dans l'ordre
surfacique--volumétrique,
\[
 N_{27}=9F_{\rm av},
 \qquad
 F_{\rm av}=
 \begin{pmatrix}0&1\\1&1\end{pmatrix}.
\]
Si
\[
 \Omega_{\rm sw}(\gamma)=\sum_{a\subset\gamma}\omega_{\rm sw}(a),
 \qquad
 S(\gamma)=\sum_{a\subset\gamma}|\omega_{\rm sw}(a)|,
 \qquad
 G_{\rm or}(\gamma)=\sum_{a\subset\gamma}g(a),
\]
alors un bloc de vingt-sept pas contient trois tours nonadiques et une
inversion d'orientation :
\[
 \Omega_{\rm sw}(\mathcal K_{27})=0,\qquad
 S(\mathcal K_{27})=18,\qquad
 G_{\rm or}(\mathcal K_{27})=1.
\]
Dans le retour observable de \(108\) pas,
\[
 \Omega_{\rm sw}(C_{108})=0,\qquad
 S(C_{108})=72,\qquad
 G_{\rm or}(C_{108})=4.
\]
Ainsi, pour la fonctionnelle additive déclarée
\[
 \mathcal A_\ast(\gamma)
 :=G_{\rm or}(\gamma)+\lambda_\ast S(\gamma),
\]
on obtient, sans introduire d'échelle dimensionnelle,
\[
 \mathcal A_\ast(\mathcal K_{27})=1+18\lambda_\ast,
 \qquad
 \mathcal A_\ast(C_{108})=4+72\lambda_\ast.
\]
Ces valeurs sont des comptages du registre. Leur utilisation thermique ou
gravitationnelle ultérieure n'altère pas leur provenance.

\subsection{Système dirigé de revêtements, raffinement et limite de la connexion de Cartan}
\label{sec:cobertura-acumulativa-dirigida}

Soit \(\mathcal O_{108}\) l'orbite observable et, pour chaque arête chronologique
\(e\), soit
\[
 c(e)=\bigl(g(e),s(e)\bigr)\in\mathbb N^2
\]
le couple formé par les indicateurs d'inversion d'orientation et de changement
effectif de feuillet. Pour une histoire dirigée
\(\gamma=e_1\cdots e_n\), on définit
\[
 c(\gamma)=\sum_{j=1}^{n}c(e_j).
\]
La loi de concaténation prend la forme
\[
 c_{m+n}(x)
 =c_m(x)+c_n\!\left(F_{\rm obs}^{m}x\right).
\]
Les comptages précédents expriment, uniformément sur l'orbite,
\[
 c_{27}(x)=(1,18),
 \qquad
 c_{108}(x)=(4,72).
\]

\begin{theorem}[Revêtement cumulatif de la chronologie observable]
\label{thm:cobertura-acumulativa-dirigida}
Sur
\[
 \widetilde{\mathcal O}_{108}^{+}
 :=\mathcal O_{108}\times\mathbb N^2
\]
définissons
\[
 \widetilde F^{+}(x,m)
 :=\bigl(F_{\rm obs}x,m+c(x\to F_{\rm obs}x)\bigr).
\]
Alors le transport de vingt-sept pas satisfait
\[
 \widetilde{\mathcal K}_{27}^{+}(x,m)
 =\bigl(\mathcal K_{27}x,m+(1,18)\bigr),
\]
et
\[
 \left(\widetilde{\mathcal K}_{27}^{+}\right)^4(x,m)
 =\bigl(x,m+(4,72)\bigr)\neq(x,m).
\]
Par conséquent, sa projection observable est d'ordre quatre, tandis que
l'action cumulative du monoïde chronologique possède un ordre de semi-groupe infini.
\end{theorem}

\begin{proof}
La loi cocyclique et la constance de \(c_{27}\) donnent
\[
 c_{108}(x)
 =\sum_{j=0}^{3}c_{27}(\mathcal K_{27}^{j}x)
 =4(1,18)=(4,72).
\]
La première coordonnée revient par
\(\mathcal K_{27}^{4}=I\) ; la seconde ne revient pas. Après \(q>0\) itérations,
l'incrément est \(q(4,72)\), non nul.
\end{proof}

\begin{remark}[Type de l'extension cumulative]
La construction précédente étend les histoires dirigées au moyen de compteurs
positifs. Ce n'est pas un relèvement automorphe du groupoïde
\(\mathcal G_{\rm mem}\), et elle ne démontre pas que la coordonnée illimitée appartienne
déjà à l'état \(\mathcal X_{\mathrm{TPK}}^{\mathrm{enr}}\). Sa complétion de Grothendieck
\(\mathbb Z^2\) peut être utilisée pour l'analyse harmonique, mais elle ne transforme pas
le coût positif
\(\mathcal A_\ast=G_{\rm or}+\lambda_\ast S\) en un caractère du groupe
d'isotropie : lorsqu'on inverse formellement une voie, ses indicateurs ne prennent pas
le signe opposé.
\end{remark}

\paragraph{Domaine de validité de l'extension cumulative.}
La hiérarchie \(27\to54\to108\), les incidences de \(F_{\rm av}\) et les
cocycles précédents sont des identités exactes dans la chronologie et les
lecteurs déclarés. L'affirmation est restreinte à l'état observable que
ces lecteurs conservent ; l'étendre à l'état enrichi complet exigerait
de prouver le retour de chaque champ oublié.

\section{Groupoïde de mémoire}
\label{sec:grupoide-memoria-tpk}

Soit \(\mathsf P_{\rm TPK}^{\rm enr}\) le groupoïde des chemins enrichis :
ses objets sont des états TPK survivants et ses flèches sont des voies admissibles
avec origine et destination déclarées, composition par concaténation et inversion
par renversement de voie. Soit \(\sim_{\rm mem}\) une équivalence qui
n'identifie deux flèches que si elle conserve les champs de mémoire requis par
le lecteur considéré et qui est congruente avec les identités, l'inversion et la
concaténation. Nous définissons
\[
 \mathcal G_{\rm mem}
 :=\mathsf P_{\rm TPK}^{\rm enr}/\!\sim_{\rm mem}.
\]
Cette définition ne transforme pas la mémoire en torsion. La mémoire appartient à
l'état discret et distingue des histoires ; la torsion appartient à une connexion
géométrique et mesure \(D_\omega e\). Seule une application entre les deux domaines peut
les relier.

Une réalisation géométrique exige, comme données supplémentaires, une longueur
\(\ell_0>0\), une représentation du groupoïde
\[
 \rho_{\rm C}:
 \mathcal G_{\rm mem}
 \longrightarrow
 \operatorname{Spin}(1,3)\ltimes\RR^{1,3},
\]
et une application de champs
\[
 \mathfrak R_{\rm grav}:
 \mathsf P_{\rm TPK}^{\rm enr}
 \dashrightarrow
 \mathsf{Cartan}_{1,3}.
\]
Pour définir la seconde application, il faut spécifier son action sur les objets et les
voies, sa compatibilité avec la concaténation et l'inversion et la régularité de la
connexion résultante. Aucune de ces propriétés ne se déduit du simple comptage
des retours.

\section{Connexion conjointe et condition d'entrelacement}
\label{sec:conexion-cartan-conjunta}

Sous les hypothèses géométriques du \cref{chap:gravedad} ---variété
orientée de dimension quatre, cotétrade \(e^I\), connexion de Lorentz
\(\omega^{IJ}\) et conditions au bord qui y sont déclarées---, la connexion de
Cartan conjointe s'écrit
\[
 \mathcal A_{\ell_0}
 =
 \begin{pmatrix}
  \omega&\ell_0^{-1}e\\
  0&0
 \end{pmatrix}.
\]
Sa courbure est l'identité algébrique
\[
 \mathcal F_{\ell_0}
 =\dd\mathcal A_{\ell_0}
  +\mathcal A_{\ell_0}\wedge\mathcal A_{\ell_0}
 =
 \begin{pmatrix}
  F_\omega&\ell_0^{-1}T_{\rm tor}\\
  0&0
 \end{pmatrix},
\]
où
\[
 F_\omega^{IJ}
 =\dd\omega^{IJ}+\omega^I{}_K\wedge\omega^{KJ},
 \qquad
 T_{\rm tor}^{I}=D_\omega e^I.
\]
La courbure et la torsion apparaissent comme des composantes distinctes d'un même défaut
de transport. L'égalité matricielle découle algébriquement des données
explicites \(e,\omega,\ell_0\), mais ne construit pas ces entrées à partir de TPK.

Le contenu fort du pont est la commutativité des holonomies :
\[
 \boxed{
 \Hol_{\mathcal A_{\ell_0}}
 \bigl(\mathfrak R_{\rm grav}(\gamma)\bigr)
 =
 \rho_{\rm C}\!\left(
 \Hol_{\rm TPK}(\gamma)
 \right)}
 \qquad
 (\gamma\in\mathsf P_{\rm TPK}^{\rm enr}).
\]
L'égalité doit être compatible avec les identités, l'inversion, la concaténation,
la subdivision, le changement de point de base et la conjugaison. Alors seulement la mémoire
discrète est représentée dans une connexion ; elle n'est pas rebaptisée
torsion ou courbure.

\begin{criterion}[Conditions d'existence du pont TPK--Cartan]
\label{crit:puente-tpk-cartan}
Il n'existe pas de réalisation possédant la commutativité et la compatibilité déclarées
si l'un quelconque des faits suivants se produit :
\begin{enumerate}
 \item deux représentants équivalents d'une même voie produisent
 des holonomies géométriques non conjuguées ;
 \item l'image ne préserve pas les identités, les inverses ou la concaténation ;
 \item une subdivision admissible altère l'holonomie ;
 \item l'égalité des holonomies échoue pour \(\mathcal K_{27}\) ou pour ses
 concaténations de profondeurs \(54\) et \(108\) ;
 \item l'application identifie mémoire et torsion sans exhiber un morphisme qui
 préserve leurs types ;
 \item un choix descendant ---l'action de Holst, \(G\), une masse ou
 une carte SI--- est utilisé pour redéfinir \(A\), \(C^\ast\) ou le retour
 TPK d'origine.
\end{enumerate}
\end{criterion}

\paragraph{Antériorité de la fonctionnelle \(\Gamma_{6\to8}\) par rapport à la réalisation Palatini--Cartan--Holst.}
Le \cref{mdg:chap:barbero} fixe, une fois \(A,C^\ast,q_\pm\) construits, la
fonctionnelle interne \(\Gamma_{6\to8}\). Le \cref{chap:gravedad} étudie ensuite
l'action Palatini--Cartan--Holst et ses équations variationnelles. Le présent pont
n'introduit pas de voie inverse entre les deux étapes.
